\chapter{Temporal Trajectories and Descent over Input Interfaces}
\label{ch:temporal-trajectories-interface-descent}

\section{Overview}
\label{sec:temporal-descent-overview}

This chapter develops two complementary locality principles for categorical
automata:
\[
    \text{locality in time}
    \qquad\text{and}\qquad
    \text{locality in the input interface}.
\]
Temporal locality says that a finite execution may be restricted to shorter
temporal windows and reconstructed by gluing compatible adjacent execution
segments. Interface locality says that a system over a global input interface
may, under suitable hypotheses, be reconstructed from compatible systems over
local interface patches.

The two forms of locality occur at different categorical levels. Temporal
restriction varies the extent of an execution while leaving the input
interface fixed. Interface restriction varies the input category itself and
therefore moves between fibres of the indexed category of Mealy systems.

The principal constructions of the chapter are:
\begin{enumerate}
    \item the discrete interval sheaf of paths of an internal graph;
    \item the Segal enhancement of this interval sheaf for an internal
    category;
    \item the temporal trajectory sheaf of the execution category of a Mealy
    morphism;
    \item the indexed category of strict Mealy systems over varying input
    interfaces;
    \item extensive and effective descent for suitable classes of input
    interface covers;
    \item descent of residual behavioural contracts and compatibility of
    behavioural abstraction with admissible input restriction;
    \item the interchange between temporal gluing and interface descent.
\end{enumerate}

The systems-theoretic interpretation is that a covering family of input
interfaces specifies a family of partial boundary views which is collectively
sufficient for reconstruction. Depending on the application, the patches may
represent sensor views, port sets, operating modes, local disturbance classes,
distributed controllers, or overlapping engineering models. Effective
descent then expresses integrability of compatible local models, while failure
of descent records a genuine global obstruction.

The chapter does not introduce a new dynamic-logic syntax. Diamonds, boxes,
tests, guarded commands, and predicate-transformer semantics are developed in
Chapter~\ref{ch:relative-mealy-program-categories}. The present chapter
supplies temporal and interface-local semantic structure which may later be
transported into that logical setting.

Unless otherwise stated, \(\mathcal E\) is finitely complete. Stronger
hypotheses are introduced locally when regular images, extensive
decompositions, dependent products, or effective descent are required.

\section{Locality as a systems principle}
\label{sec:locality-systems-principle}

\subsection{Temporal and interface locality}
\label{subsec:temporal-interface-locality}

Let
\[
    \Phi:\mathbb A\rightsquigarrow\mathbb B
\]
be a Mealy morphism.

Temporal locality concerns finite executions of \(\Phi\). A long execution
contains shorter contiguous subexecutions, and compatible adjacent
subexecutions may be concatenated.

Interface locality concerns input functors
\[
    u:\mathbb U\to\mathbb A.
\]
The functor \(u\) is regarded as a local view, patch, refinement, or
presentation of the global input interface \(\mathbb A\). Pullback of
\(\Phi\) along \(u\) gives a restricted system
\[
    u^\ast\Phi:\mathbb U\rightsquigarrow\mathbb B.
\]

These two restrictions are independent:
\begin{enumerate}
    \item temporal restriction changes the observed execution interval;
    \item input restriction changes the available input interface.
\end{enumerate}

\subsection{Identifiability, integrability, and effectivity}
\label{subsec:identifiability-integrability-effectivity}

Let
\[
    \{u_i:\mathbb U_i\to\mathbb A\}_{i\in I}
\]
be a covering family of input interfaces.

\begin{definition}[Interface identifiability]
\label{def:interface-identifiability}
The covering family is \textbf{identifying} for a class of systems if, whenever
\[
    \Phi,\Psi:\mathbb A\rightsquigarrow\mathbb B
\]
have isomorphic restrictions
\[
    u_i^\ast\Phi\cong u_i^\ast\Psi
\]
compatible on all overlaps, then
\[
    \Phi\cong\Psi.
\]
\end{definition}

\begin{definition}[Interface integrability]
\label{def:interface-integrability}
The covering family is \textbf{integrating} for a class of systems if every
compatible family of local systems
\[
    \Phi_i:\mathbb U_i\rightsquigarrow\mathbb B
\]
with coherent overlap identifications arises from a global system
\[
    \Phi:\mathbb A\rightsquigarrow\mathbb B.
\]
\end{definition}

\begin{definition}[Effective interface descent]
\label{def:effective-interface-descent-informal}
A covering family has \textbf{effective descent} for a class of systems if it
is both identifying and integrating, with the reconstruction unique up to
unique coherent isomorphism.
\end{definition}

Identifiability is an observational condition: the local views jointly
distinguish global systems. Integrability is a synthesis condition:
compatible local implementations may be assembled globally. Effectivity says
that the local-to-global process captures precisely the global systems.

\subsection{Systems-theoretic interpretations}
\label{subsec:systems-interpretations-descent}

A family of input patches may represent:
\begin{enumerate}
    \item different sensors observing overlapping variables;
    \item local port sets in a networked architecture;
    \item subsystem boundaries with shared boundary variables;
    \item operating-mode fragments;
    \item local classes of disturbances or control inputs;
    \item information available to distributed controllers;
    \item overlapping physical, software, and control-theoretic models.
\end{enumerate}

Compatibility on overlaps means that the local systems agree on shared input
behaviour, shared state identifications, and shared transition-output data.
The cocycle equation expresses higher coherence among these identifications.

Failure of descent may therefore indicate:
\begin{enumerate}
    \item an omitted global degree of freedom;
    \item incompatible local control laws;
    \item a hidden coupling invisible on individual patches;
    \item pairwise consistency without triple consistency;
    \item a locally valid but globally unrealizable specification.
\end{enumerate}

\begin{remark}[Ambient sheaf semantics]
\label{rem:ambient-sheaf-semantics}
When
\[
    \mathcal E\simeq\operatorname{Sh}(\mathcal C,J)
\]
is presented as a Grothendieck topos, all objects, images, quotients, and
subobjects used below already admit their usual local interpretation over the
ambient site. This observation is not the principal sheaf-theoretic content
of the chapter. The new structure concerns sheaves in a temporal variable and
descent in the input-interface variable.
\end{remark}

\section{Discrete interval sheaves}
\label{sec:discrete-interval-sheaves}

\subsection{The discrete interval category}
\label{subsec:discrete-interval-category}

\begin{definition}[Discrete interval category]
\label{def:discrete-interval-category}
The \textbf{discrete interval category}
\[
    \mathbf{Int}_{\mathbb N}
\]
has the natural numbers as objects. A morphism
\[
    (i,k):k\longrightarrow n
\]
is an inclusion of a contiguous interval of length \(k\) into one of length
\(n\), determined by an integer \(i\) satisfying
\[
    0\leq i\leq n-k.
\]
Composition is addition of starting positions:
\[
    (i,k):k\to n,
    \qquad
    (j,\ell):\ell\to k
\]
compose to
\[
    (i+j,\ell):\ell\to n.
\]
\end{definition}

An object \(n\) is regarded as a temporal extent containing \(n\) primitive
steps and \(n+1\) boundary instants. A morphism selects a contiguous
subinterval.

\begin{definition}[Internal interval presheaf]
\label{def:internal-interval-presheaf}
An \(\mathcal E\)-valued \textbf{interval presheaf} is a functor
\[
    X:\mathbf{Int}_{\mathbb N}^{\mathrm{op}}\to\mathcal E.
\]
For a subinterval inclusion
\[
    \iota:k\to n,
\]
the induced map
\[
    X(\iota):X(n)\to X(k)
\]
is called temporal restriction.
\end{definition}

\subsection{Adjacent covers}
\label{subsec:adjacent-interval-covers}

For \(m,n\in\mathbb N\), the interval \(m+n\) is covered by its initial
subinterval of length \(m\) and its terminal subinterval of length \(n\).
Their intersection is the common boundary point.

The corresponding cover is represented by
\[
    m\longrightarrow m+n
    \longleftarrow n.
\]

\begin{definition}[Discrete interval sheaf]
\label{def:discrete-interval-sheaf}
An interval presheaf
\[
    X:\mathbf{Int}_{\mathbb N}^{\mathrm{op}}\to\mathcal E
\]
is a \textbf{discrete interval sheaf} if, for every \(m,n\), the canonical map
\[
    X(m+n)\longrightarrow
    X(m)\times_{X(0)}X(n)
\]
is an isomorphism, where the maps to \(X(0)\) are restriction to the common
boundary point.
\end{definition}

Thus compatible adjacent sections glue uniquely.

\subsection{Trajectory sheaves of internal graphs}
\label{subsec:trajectory-sheaves-internal-graphs}

Let
\[
    G=
    \left(
    G_1
    \mathrel{\substack{\xrightarrow{t_G}\\[-0.4em]\xrightarrow[s_G]{}}}
    G_0
    \right)
\]
be an internal graph in \(\mathcal E\).

\begin{definition}[Path objects]
\label{def:internal-graph-path-objects}
Define
\[
    G^{[0]}:=G_0,
    \qquad
    G^{[1]}:=G_1,
\]
and, for \(n\geq 2\),
\[
    G^{[n]}
    :=
    \underbrace{
    G_1\times_{G_0}G_1\times_{G_0}\cdots\times_{G_0}G_1
    }_{n\text{ factors}},
\]
where each pullback identifies the target of one edge with the source of the
next edge.
\end{definition}

A generalized element of \(G^{[n]}\) is a composable graph path
\[
    x_0\xrightarrow{g_1}x_1
    \xrightarrow{g_2}\cdots
    \xrightarrow{g_n}x_n.
\]

\begin{definition}[Trajectory presheaf]
\label{def:trajectory-presheaf-internal-graph}
The \textbf{trajectory presheaf} of \(G\) is
\[
    \operatorname{Traj}(G):
    \mathbf{Int}_{\mathbb N}^{\mathrm{op}}\to\mathcal E
\]
defined by
\[
    \operatorname{Traj}(G)(n):=G^{[n]}.
\]
Restriction to a subinterval selects the corresponding contiguous subpath.
\end{definition}

\begin{theorem}[Temporal gluing of graph trajectories]
\label{thm:temporal-gluing-graph-trajectories}
For every internal graph \(G\), the trajectory presheaf
\[
    \operatorname{Traj}(G)
\]
is a discrete interval sheaf.
\end{theorem}

\begin{proof}
For \(m,n\in\mathbb N\), a path of length \(m+n\) determines:
\begin{enumerate}
    \item its initial path of length \(m\);
    \item its terminal path of length \(n\);
    \item a common boundary object at the endpoint of the first path and the
    starting point of the second.
\end{enumerate}
This gives a canonical map
\[
    G^{[m+n]}
    \longrightarrow
    G^{[m]}\times_{G_0}G^{[n]}.
\]
Conversely, a pair of paths with equal common boundary object determines a
unique concatenated edge string of length \(m+n\). Since all path objects are
defined by iterated pullbacks, the two constructions are inverse by the
universal property of pullbacks.
\end{proof}

\subsection{Reconstruction of the underlying graph}
\label{subsec:reconstruction-underlying-graph}

Let \(X\) be a discrete interval sheaf. Define
\[
    G_X{}_0:=X(0),
    \qquad
    G_X{}_1:=X(1).
\]
The two endpoint inclusions
\[
    0\to 1
\]
induce source and target maps
\[
    s_X,t_X:X(1)\to X(0).
\]

\begin{proposition}[Reconstruction from lengths zero and one]
\label{prop:reconstruction-length-zero-one}
For every discrete interval sheaf \(X\), there are canonical isomorphisms
\[
    X(n)
    \cong
    \underbrace{
    X(1)\times_{X(0)}\cdots\times_{X(0)}X(1)
    }_{n\text{ factors}}.
\]
Hence \(X\) is canonically isomorphic to
\[
    \operatorname{Traj}(G_X).
\]
\end{proposition}

\begin{proof}
Proceed by induction on \(n\). The cases \(n=0,1\) are immediate. For
\(n+1\), apply the sheaf condition to the decomposition into an interval of
length \(n\) followed by one of length \(1\):
\[
    X(n+1)\cong X(n)\times_{X(0)}X(1).
\]
The inductive hypothesis gives the required iterated pullback.
\end{proof}

\begin{theorem}[Internal graph--interval-sheaf equivalence]
\label{thm:internal-graph-interval-sheaf-equivalence}
The trajectory construction defines an equivalence
\[
    \operatorname{Traj}:
    \mathbf{Graph}_{\mathrm{int}}(\mathcal E)
    \simeq
    \mathbf{Sh}_{\mathcal E}(\mathbf{Int}_{\mathbb N})
\]
between internal graphs in \(\mathcal E\) and
\(\mathcal E\)-valued discrete interval sheaves.
\end{theorem}

\begin{proof}
The preceding proposition constructs a quasi-inverse by sending an interval
sheaf \(X\) to the internal graph
\[
    X(1)
    \mathrel{\substack{\xrightarrow{t_X}\\[-0.4em]\xrightarrow[s_X]{}}}
    X(0).
\]
The unit and counit are the canonical path-reconstruction isomorphisms.
\end{proof}

\section{Segal enhancement and categorical execution}
\label{sec:segal-enhancement-categorical-execution}

\subsection{The nerve of an internal category}
\label{subsec:nerve-internal-category}

Let
\[
    \mathbb C=(C_0,C_1,s_C,t_C,e_C,m_C)
\]
be an internal category.

\begin{definition}[Internal nerve]
\label{def:internal-nerve}
The \textbf{nerve} of \(\mathbb C\) is the simplicial object
\[
    N\mathbb C:\Delta^{\mathrm{op}}\to\mathcal E
\]
with
\[
    N\mathbb C_0=C_0
\]
and
\[
    N\mathbb C_n
    =
    \underbrace{
    C_1\times_{C_0}\cdots\times_{C_0}C_1
    }_{n\text{ factors}}.
\]
The outer face maps delete the initial or terminal arrow, the inner face maps
compose adjacent arrows, and the degeneracy maps insert identity arrows.
\end{definition}

\subsection{Segal identities}
\label{subsec:segal-identities}

\begin{proposition}[Segal isomorphisms]
\label{prop:internal-category-segal-isomorphisms}
For every \(n\geq 2\), the canonical Segal map
\[
    N\mathbb C_n
    \longrightarrow
    \underbrace{
    N\mathbb C_1
    \times_{N\mathbb C_0}\cdots
    \times_{N\mathbb C_0}
    N\mathbb C_1
    }_{n\text{ factors}}
\]
is an isomorphism.
\end{proposition}

\begin{proof}
Both sides are the same iterated pullback of \(C_1\) over \(C_0\).
\end{proof}

\subsection{Temporal shadow}
\label{subsec:temporal-shadow-nerve}

Let \(U\mathbb C\) denote the underlying internal graph of \(\mathbb C\).

\begin{definition}[Temporal shadow]
\label{def:temporal-shadow}
The \textbf{temporal shadow} of \(\mathbb C\) is the interval sheaf
\[
    \operatorname{Temp}(\mathbb C)
    :=
    \operatorname{Traj}(U\mathbb C).
\]
\end{definition}

The objects of
\[
    \operatorname{Temp}(\mathbb C)(n)
\]
and
\[
    N\mathbb C_n
\]
coincide. Their difference lies in the available structure maps. The
interval sheaf records restriction to contiguous subpaths. The nerve
additionally records identity insertion and composition of adjacent arrows.

\begin{remark}[Information retained by the Segal enhancement]
\label{rem:information-segal-enhancement}
The interval sheaf retains:
\begin{enumerate}
    \item objects;
    \item primitive arrows;
    \item finite composable strings;
    \item temporal restriction;
    \item temporal concatenation along matching boundaries.
\end{enumerate}
The full nerve additionally retains:
\begin{enumerate}
    \item identity insertion;
    \item contraction of adjacent steps by categorical composition;
    \item associativity coherence among contractions.
\end{enumerate}
\end{remark}

\begin{theorem}[Segal enhancement theorem]
\label{thm:segal-enhancement}
Every internal category \(\mathbb C\) determines:
\begin{enumerate}
    \item a discrete interval sheaf
    \[
        \operatorname{Temp}(\mathbb C);
    \]
    \item a canonical Segal enhancement
    \[
        N\mathbb C.
    \]
\end{enumerate}
Every internal functor
\[
    F:\mathbb C\to\mathbb D
\]
induces a morphism of interval sheaves
\[
    \operatorname{Temp}(F):
    \operatorname{Temp}(\mathbb C)
    \to
    \operatorname{Temp}(\mathbb D)
\]
and a simplicial map
\[
    NF:N\mathbb C\to N\mathbb D.
\]
\end{theorem}

\begin{proof}
The nerve is functorial in internal functors. Restricting to the maps which
select contiguous subpaths gives the interval-sheaf morphism.
\end{proof}

\section{Temporal trajectories of Mealy morphisms}
\label{sec:temporal-trajectories-mealy}

\subsection{Execution categories}
\label{subsec:execution-categories-temporal}

Let
\[
    \Phi=(P,\delta_P,\omega_P):
    \mathbb A\rightsquigarrow\mathbb B
\]
be a Mealy morphism.

Recall that the carrier \(P\) carries a right \(\mathbb A\)-action by
\[
    \delta_P:
    A_1\times_{t_A,A_0,p_P}P\to P.
\]
Its action category is
\[
    \int_{\mathbb A}P.
\]
Explicitly,
\[
    \left(\int_{\mathbb A}P\right)_0=P
\]
and
\[
    \left(\int_{\mathbb A}P\right)_1
    =
    A_1\times_{t_A,A_0,p_P}P.
\]
The source and target maps are
\[
    s(a,x)=x,
    \qquad
    t(a,x)=\delta_P(a,x).
\]

\subsection{Input and output labelling}
\label{subsec:input-output-labelling-temporal}

The input projection is the internal functor
\[
    I_P:
    \int_{\mathbb A}P\to\mathbb A^{\mathrm{op}}
\]
with object part
\[
    p_P:P\to A_0
\]
and arrow part
\[
    (a,x)\longmapsto a.
\]

The output structure determines an internal functor
\[
    O_P:
    \int_{\mathbb A}P\to\mathbb B^{\mathrm{op}}
\]
with object part
\[
    q_P:P\to B_0
\]
and arrow part
\[
    (a,x)\longmapsto\omega_P(a,x),
\]
viewed oppositely.

The direction into \(\mathbb B^{\mathrm{op}}\) agrees with the convention that
\[
    \omega_P(a,x):
    q_P(\delta_P(a,x))\to q_P(x)
\]
is an arrow in \(\mathbb B\).

\subsection{Temporal execution sheaf}
\label{subsec:temporal-execution-sheaf}

\begin{definition}[Temporal execution sheaf]
\label{def:temporal-execution-sheaf}
The \textbf{temporal execution sheaf} of \(\Phi\) is
\[
    \mathscr T_\Phi
    :=
    \operatorname{Temp}
    \left(
    \int_{\mathbb A}P
    \right).
\]
\end{definition}

Its object of \(n\)-step executions is
\[
    \mathscr T_\Phi(n)
    =
    \underbrace{
    \left(A_1\times_{A_0}P\right)
    \times_P\cdots\times_P
    \left(A_1\times_{A_0}P\right)
    }_{n\text{ factors}}.
\]

In generalized-element notation, an \(n\)-step execution has the form
\[
    x_0
    \xrightarrow{a_1}
    x_1
    \xrightarrow{a_2}
    \cdots
    \xrightarrow{a_n}
    x_n,
\]
where
\[
    x_i=\delta_P(a_i,x_{i-1})
\]
with the target-to-source order determined by the standing conventions.

The input and output functors induce a span of interval sheaves
\[
\begin{tikzcd}[column sep=large]
    \operatorname{Temp}(\mathbb A^{\mathrm{op}})
    &
    \mathscr T_\Phi
        \arrow[l, "\operatorname{Temp}(I_P)"']
        \arrow[r, "\operatorname{Temp}(O_P)"]
    &
    \operatorname{Temp}(\mathbb B^{\mathrm{op}}).
\end{tikzcd}
\]

\subsection{Segal-enhanced temporal semantics}
\label{subsec:segal-enhanced-temporal-semantics}

\begin{definition}[Segal-enhanced execution semantics]
\label{def:segal-enhanced-execution-semantics}
The \textbf{Segal-enhanced temporal semantics} of \(\Phi\) is the span
\[
\begin{tikzcd}[column sep=large]
    N(\mathbb A^{\mathrm{op}})
    &
    N\left(\int_{\mathbb A}P\right)
        \arrow[l, "NI_P"']
        \arrow[r, "NO_P"]
    &
    N(\mathbb B^{\mathrm{op}}).
\end{tikzcd}
\]
\end{definition}

The interval-sheaf span is obtained by forgetting the inner face and
degeneracy maps while retaining restriction to contiguous temporal
subintervals.

\subsection{Functoriality in strict semantic homomorphisms}
\label{subsec:temporal-functoriality-strict}

Let
\[
    f:\Phi\to\Psi
\]
be a strict semantic homomorphism between Mealy morphisms
\[
    \Phi=(P,\delta_P,\omega_P),
    \qquad
    \Psi=(Q,\delta_Q,\omega_Q).
\]

The strict preservation equations induce an internal functor
\[
    \int_{\mathbb A}f:
    \int_{\mathbb A}P
    \to
    \int_{\mathbb A}Q
\]
with object part \(f\) and arrow part
\[
    (a,x)\longmapsto(a,f(x)).
\]

\begin{proposition}[Functoriality of temporal execution]
\label{prop:functoriality-temporal-execution}
Every strict semantic homomorphism
\[
    f:\Phi\to\Psi
\]
induces:
\begin{enumerate}
    \item a morphism of temporal execution sheaves
    \[
        \mathscr T_f:
        \mathscr T_\Phi\to\mathscr T_\Psi;
    \]
    \item a simplicial map
    \[
        N\left(\int_{\mathbb A}f\right):
        N\left(\int_{\mathbb A}P\right)
        \to
        N\left(\int_{\mathbb A}Q\right).
    \]
\end{enumerate}
Both maps commute with the input and output labelling.
\end{proposition}

\begin{proof}
Apply Theorem~\ref{thm:segal-enhancement} to the internal functor
\[
    \int_{\mathbb A}f.
\]
Strict preservation of input and output gives commutativity with the labelled
spans.
\end{proof}

\subsection{Temporal semantics and Mealy composition}
\label{subsec:temporal-semantics-mealy-composition}

Let
\[
    \Phi:\mathbb A\rightsquigarrow\mathbb B,
    \qquad
    \Psi:\mathbb B\rightsquigarrow\mathbb C
\]
be composable Mealy morphisms.

An \(n\)-step execution of the composite
\[
    \Psi\circ\Phi
\]
consists of:
\begin{enumerate}
    \item an \(n\)-step execution of \(\Phi\);
    \item an \(n\)-step execution of \(\Psi\);
    \item equality between the output trajectory of \(\Phi\) and the input
    trajectory of \(\Psi\).
\end{enumerate}

\begin{theorem}[Temporal composition]
\label{thm:temporal-composition}
For every \(n\in\mathbb N\), there is a canonical isomorphism
\[
    \mathscr T_{\Psi\circ\Phi}(n)
    \cong
    \mathscr T_\Phi(n)
    \times_{\operatorname{Temp}(\mathbb B^{\mathrm{op}})(n)}
    \mathscr T_\Psi(n).
\]
These isomorphisms are natural in interval restrictions and therefore assemble
into an isomorphism of interval sheaves
\[
    \mathscr T_{\Psi\circ\Phi}
    \cong
    \mathscr T_\Phi
    \times_{\operatorname{Temp}(\mathbb B^{\mathrm{op}})}
    \mathscr T_\Psi.
\]
\end{theorem}

\begin{proof}
For \(n=0\), the assertion is the pullback definition of the composite
carrier.

For \(n=1\), the assertion is the pullback description of a primitive
execution of the composite: a primitive execution in \(\Phi\) and one in
\(\Psi\) whose intermediate \(\mathbb B\)-arrow agrees.

For general \(n\), use the Segal decomposition
\[
    \mathscr T_\Phi(n)
    \cong
    \underbrace{
    \mathscr T_\Phi(1)
    \times_{\mathscr T_\Phi(0)}
    \cdots
    \times_{\mathscr T_\Phi(0)}
    \mathscr T_\Phi(1)
    }_{n\text{ factors}},
\]
and similarly for \(\Psi\) and \(\Psi\circ\Phi\). Finite limits commute with
finite limits, so the length-one pullback comparison extends canonically to
all lengths. Naturality under temporal restriction follows from the
construction.
\end{proof}

\begin{remark}[Temporal abstract machine]
\label{rem:temporal-abstract-machine}
The span
\[
    \operatorname{Temp}(\mathbb A^{\mathrm{op}})
    \longleftarrow
    \mathscr T_\Phi
    \longrightarrow
    \operatorname{Temp}(\mathbb B^{\mathrm{op}})
\]
is an interval-sheaf abstract machine derived from the stateful Mealy
structure. The Segal enhancement records additional compositional structure
which is absent from the underlying interval sheaf alone.
\end{remark}

\section{Systems indexed by input interfaces}
\label{sec:systems-indexed-input-interfaces}

\subsection{Fixing the output interface}
\label{subsec:fixing-output-interface}

Fix an internal category \(\mathbb B\).

The output interface remains fixed throughout the descent theory below. The
input variable is contravariant under input pullback and therefore gives the
natural base of an indexed category.

\subsection{Input pullback}
\label{subsec:input-pullback-systems}

Let
\[
    F:\mathbb A'\to\mathbb A
\]
be an internal functor, and let
\[
    \Phi=(P,\delta_P,\omega_P):
    \mathbb A\rightsquigarrow\mathbb B.
\]

The input pullback
\[
    F^\ast\Phi:
    \mathbb A'\rightsquigarrow\mathbb B
\]
has carrier
\[
    P_F
    :=
    A'_0\times_{F_0,A_0,p_P}P.
\]

A generalized element is written
\[
    (v',x),
    \qquad
    F_0(v')=p_P(x).
\]

For an arrow
\[
    a':u'\to v'
\]
of \(\mathbb A'\), define
\[
    \delta_{P_F}(a',(v',x))
    :=
    \left(
    u',
    \delta_P(F_1a',x)
    \right),
\]
and
\[
    \omega_{P_F}(a',(v',x))
    :=
    \omega_P(F_1a',x).
\]

\begin{proposition}[Input pullback]
\label{prop:input-pullback-systems}
The displayed formulas define a Mealy morphism
\[
    F^\ast\Phi:
    \mathbb A'\rightsquigarrow\mathbb B.
\]
Input pullback extends to a functor
\[
    F^\ast:
    \mathsf{Mly}^{\mathrm{str}}_{\mathbb A,\mathbb B}
    \to
    \mathsf{Mly}^{\mathrm{str}}_{\mathbb A',\mathbb B}.
\]
\end{proposition}

\begin{proof}
The typing, unit, and multiplication equations follow by applying the
corresponding equations of \(\Phi\) to the arrows \(F_1a'\). Strict maps are
pulled back on carriers, and their transition-output preservation equations
remain valid after pullback.
\end{proof}

\subsection{The indexed category of systems}
\label{subsec:indexed-category-systems}

\begin{definition}[Input-indexed system category]
\label{def:input-indexed-system-category}
Define
\[
    \mathsf{Sys}_{\mathbb B}:
    \mathsf{IntCat}(\mathcal E)^{\mathrm{op}}
    \to
    \mathbf{Cat}
\]
by
\[
    \mathsf{Sys}_{\mathbb B}(\mathbb A)
    :=
    \mathsf{Mly}^{\mathrm{str}}_{\mathbb A,\mathbb B}
\]
and
\[
    \mathsf{Sys}_{\mathbb B}(F):=F^\ast.
\]
\end{definition}

Input pullback is pseudofunctorial:
\[
    1_{\mathbb A}^\ast\cong 1,
\]
and
\[
    (FG)^\ast\cong G^\ast F^\ast.
\]
These isomorphisms arise from canonical associativity isomorphisms of
pullbacks.

\subsection{Grothendieck construction}
\label{subsec:grothendieck-construction-systems}

\begin{definition}[Total category of systems]
\label{def:total-category-systems}
The Grothendieck construction
\[
    \int\mathsf{Sys}_{\mathbb B}
\]
has:
\begin{enumerate}
    \item objects \((\mathbb A,\Phi)\), where
    \[
        \Phi:\mathbb A\rightsquigarrow\mathbb B;
    \]
    \item morphisms
    \[
        (\mathbb A',\Phi')
        \longrightarrow
        (\mathbb A,\Phi)
    \]
    consisting of an internal functor
    \[
        F:\mathbb A'\to\mathbb A
    \]
    and a strict semantic homomorphism
    \[
        \Phi'\to F^\ast\Phi.
    \]
\end{enumerate}
\end{definition}

The projection
\[
    \int\mathsf{Sys}_{\mathbb B}
    \to
    \mathsf{IntCat}(\mathcal E)
\]
is the fibration of strict Mealy systems over varying input interfaces.

\subsection{Indexed behavioural predicates and contracts}
\label{subsec:indexed-predicates-contracts}

For each \(\mathbb A\), write
\[
    \mathbf{Beh}_{\mathbb A,\mathbb B,0}
\]
for the canonical behaviour carrier over
\[
    A_0\times B_0.
\]

\begin{definition}[Indexed behavioural predicates]
\label{def:indexed-behavioural-predicates}
Define
\[
    \mathsf{Pred}_{\mathbb B}(\mathbb A)
    :=
    \operatorname{Sub}_{\mathcal E/(A_0\times B_0)}
    \left(
    \mathbf{Beh}_{\mathbb A,\mathbb B,0}
    \right).
\]
\end{definition}

\begin{definition}[Indexed residual contracts]
\label{def:indexed-residual-contracts}
Define
\[
    \mathsf{Ctr}_{\mathbb B}(\mathbb A)
\]
to be the poset of derivative-invariant subobjects of
\[
    \mathbf{Beh}_{\mathbb A,\mathbb B,0}.
\]
Equivalently,
\[
    \mathsf{Ctr}_{\mathbb B}(\mathbb A)
    =
    \mathsf{Coeq}_{\mathbb A,\mathbb B}.
\]
\end{definition}

Whenever restriction of residual behaviours along
\[
    F:\mathbb A'\to\mathbb A
\]
has been fixed, it induces reindexing maps
\[
    F^\ast:
    \mathsf{Pred}_{\mathbb B}(\mathbb A)
    \to
    \mathsf{Pred}_{\mathbb B}(\mathbb A')
\]
and, under derivative compatibility,
\[
    F^\ast:
    \mathsf{Ctr}_{\mathbb B}(\mathbb A)
    \to
    \mathsf{Ctr}_{\mathbb B}(\mathbb A').
\]

\section{Extensive descent}
\label{sec:extensive-descent}

For this section, assume that \(\mathcal E\) is lextensive and admits the
small coproducts under consideration.

\subsection{Coproducts of input interfaces}
\label{subsec:coproducts-input-interfaces}

Let
\[
    \mathbb A
    =
    \coprod_{i\in I}\mathbb A_i.
\]
Then
\[
    A_0=\coprod_i A_{i,0},
    \qquad
    A_1=\coprod_i A_{i,1},
\]
with all internal-category structure maps defined componentwise.

There are no arrows connecting distinct coproduct components.

\subsection{Decomposition of carriers}
\label{subsec:decomposition-carriers}

Extensivity gives
\[
    \mathcal E/
    \left(
    \left(\coprod_iA_{i,0}\right)\times B_0
    \right)
    \simeq
    \prod_i
    \mathcal E/(A_{i,0}\times B_0).
\]

Thus a carrier
\[
    P\to
    \left(\coprod_iA_{i,0}\right)\times B_0
\]
is equivalently a family of carriers
\[
    P_i\to A_{i,0}\times B_0.
\]

\subsection{Decomposition of Mealy structure}
\label{subsec:decomposition-mealy-structure}

Since
\[
    A_1=\coprod_iA_{i,1},
\]
the derivative domain decomposes as
\[
    A_1\times_{A_0}P
    \cong
    \coprod_i
    \left(
    A_{i,1}\times_{A_{i,0}}P_i
    \right).
\]

Consequently:
\begin{enumerate}
    \item a transition map on \(P\) is equivalent to a family of transition
    maps on the \(P_i\);
    \item an output map on \(P\) is equivalent to a family of output maps on
    the \(P_i\);
    \item the Mealy unit and multiplication equations hold globally exactly
    when they hold componentwise.
\end{enumerate}

\begin{theorem}[Extensive descent for strict Mealy systems]
\label{thm:extensive-descent-systems}
There is a canonical equivalence
\[
    \mathsf{Sys}_{\mathbb B}
    \left(
    \coprod_{i\in I}\mathbb A_i
    \right)
    \simeq
    \prod_{i\in I}
    \mathsf{Sys}_{\mathbb B}(\mathbb A_i).
\]
\end{theorem}

\begin{proof}
The equivalence on carriers is supplied by extensivity. The transition and
output maps decompose over the coproduct of derivative domains, and all Mealy
equations are componentwise. Strict semantic homomorphisms decompose in the
same way.
\end{proof}

\subsection{Extensive descent for temporal trajectories}
\label{subsec:extensive-descent-trajectories}

\begin{proposition}[Temporal trajectories of coproduct systems]
\label{prop:temporal-trajectories-coproduct}
If
\[
    \Phi
    \cong
    \coprod_i\Phi_i
\]
under Theorem~\ref{thm:extensive-descent-systems}, then
\[
    \mathscr T_\Phi
    \cong
    \coprod_i\mathscr T_{\Phi_i}
\]
as interval sheaves.
\end{proposition}

\begin{proof}
The execution category decomposes as
\[
    \int_{\coprod_i\mathbb A_i}
    \left(\coprod_iP_i\right)
    \cong
    \coprod_i
    \int_{\mathbb A_i}P_i.
\]
There are no execution arrows between distinct components. Taking path
objects therefore preserves the decomposition.
\end{proof}

\subsection{Extensive descent for residual contracts}
\label{subsec:extensive-descent-contracts}

Residual behaviours over a coproduct input interface lie over one component
of the object object. Their derivatives remain in that component.

\begin{theorem}[Extensive descent for residual contracts]
\label{thm:extensive-descent-contracts}
There is a canonical isomorphism of posets
\[
    \mathsf{Ctr}_{\mathbb B}
    \left(
    \coprod_i\mathbb A_i
    \right)
    \cong
    \prod_i
    \mathsf{Ctr}_{\mathbb B}(\mathbb A_i).
\]
\end{theorem}

\begin{proof}
A subobject of the canonical behaviour carrier decomposes componentwise by
extensivity. Derivative invariance is componentwise because every input arrow
belongs to a unique coproduct component.
\end{proof}

\begin{remark}[Systems interpretation]
\label{rem:extensive-descent-systems-interpretation}
A system on a disconnected input interface is exactly a family of systems on
its connected components. No compatibility data are required because there
are no nontrivial overlaps and no input arrows connecting distinct
components.
\end{remark}

\section{Effective interface covers and descent data}
\label{sec:effective-interface-covers-descent-data}

\subsection{Singleton covers}
\label{subsec:singleton-interface-covers}

A general covering family
\[
    \{u_i:\mathbb U_i\to\mathbb A\}_{i\in I}
\]
may be replaced, when the coproduct exists, by the singleton map
\[
    p:
    \coprod_i\mathbb U_i
    \to
    \mathbb A.
\]

We therefore formulate descent primarily for a single internal functor
\[
    p:\mathbb U\to\mathbb A.
\]

\subsection{\v{C}ech nerve}
\label{subsec:cech-nerve-interface-cover}

Assume the relevant pullbacks exist. Define
\[
    \mathbb U^{[2]}
    :=
    \mathbb U\times_{\mathbb A}\mathbb U,
\]
\[
    \mathbb U^{[3]}
    :=
    \mathbb U\times_{\mathbb A}
    \mathbb U\times_{\mathbb A}\mathbb U,
\]
and similarly for higher iterated pullbacks.

Write
\[
    \pi_1,\pi_2:
    \mathbb U^{[2]}\rightrightarrows\mathbb U
\]
for the two projections, and
\[
    \pi_{12},\pi_{23},\pi_{13}:
    \mathbb U^{[3]}\to\mathbb U^{[2]}
\]
for the three projections.

\subsection{System descent data}
\label{subsec:system-descent-data}

\begin{definition}[System descent datum]
\label{def:system-descent-datum}
A \textbf{system descent datum} for
\[
    p:\mathbb U\to\mathbb A
\]
consists of:
\begin{enumerate}
    \item a local system
    \[
        \Phi_{\mathbb U}
        \in
        \mathsf{Sys}_{\mathbb B}(\mathbb U);
    \]
    \item an isomorphism
    \[
        \vartheta:
        \pi_1^\ast\Phi_{\mathbb U}
        \xrightarrow{\cong}
        \pi_2^\ast\Phi_{\mathbb U}
    \]
    in
    \[
        \mathsf{Sys}_{\mathbb B}(\mathbb U^{[2]});
    \]
    \item the cocycle equation
    \[
        \pi_{23}^\ast\vartheta
        \circ
        \pi_{12}^\ast\vartheta
        =
        \pi_{13}^\ast\vartheta
    \]
    over \(\mathbb U^{[3]}\).
\end{enumerate}
\end{definition}

Morphisms of descent data are strict semantic homomorphisms commuting with the
overlap isomorphisms.

Write
\[
    \mathsf{Desc}_p(\mathsf{Sys}_{\mathbb B})
\]
for the resulting category.

\subsection{Comparison functor}
\label{subsec:comparison-functor-system-descent}

Every global system
\[
    \Phi\in\mathsf{Sys}_{\mathbb B}(\mathbb A)
\]
gives a canonical descent datum by pullback:
\[
    p^\ast\Phi
\]
with the canonical comparison between its two pullbacks to
\(\mathbb U^{[2]}\).

\begin{definition}[System descent comparison]
\label{def:system-descent-comparison}
The \textbf{descent comparison functor} is
\[
    K_p:
    \mathsf{Sys}_{\mathbb B}(\mathbb A)
    \to
    \mathsf{Desc}_p(\mathsf{Sys}_{\mathbb B}).
\]
\end{definition}

\subsection{Component form of system descent}
\label{subsec:component-form-system-descent}

Let
\[
    \Phi_{\mathbb U}
    =
    (P_{\mathbb U},
    \delta_{\mathbb U},
    \omega_{\mathbb U}).
\]

The descent datum contains:
\begin{enumerate}
    \item descent data for the carrier
    \[
        P_{\mathbb U}
        \to
        U_0\times B_0;
    \]
    \item compatibility of the transition maps on
    \[
        U_1\times_{U_0}P_{\mathbb U};
    \]
    \item compatibility of the output maps into \(B_1\);
    \item compatibility with the carrier anchors;
    \item the local Mealy unit and multiplication laws.
\end{enumerate}

The overlap isomorphism is required to preserve transition and output
strictly.

\subsection{Contract descent data}
\label{subsec:contract-descent-data}

Suppose restriction of canonical behaviours along \(p\) has been defined.

\begin{definition}[Contract descent datum]
\label{def:contract-descent-datum}
A \textbf{contract descent datum} consists of a local residual contract
\[
    V_{\mathbb U}
    \hookrightarrow
    \mathbf{Beh}_{\mathbb U,\mathbb B,0}
\]
whose two restrictions to \(\mathbb U^{[2]}\) agree under the canonical
comparison.
\end{definition}

Since contracts form a poset, no additional nontrivial cocycle isomorphism is
required beyond equality after the canonical coherence identifications.

\section{Admissible effective covers}
\label{sec:admissible-effective-covers}

The strongest descent theorem requires a class of covers for which the
underlying categorical data, the relevant pullback objects, and the Mealy
equations descend effectively.

\subsection{Nerve components of an interface functor}
\label{subsec:nerve-components-interface-functor}

Let
\[
    p:\mathbb U\to\mathbb A
\]
be an internal functor.

For each \(n\geq 0\), the nerve induces a map
\[
    p_n:
    U_n\to A_n,
\]
where
\[
    U_n=N\mathbb U_n,
    \qquad
    A_n=N\mathbb A_n.
\]

Thus:
\[
    p_0:U_0\to A_0,
\]
\[
    p_1:U_1\to A_1,
\]
\[
    p_2:
    U_1\times_{U_0}U_1
    \to
    A_1\times_{A_0}A_1,
\]
and similarly in higher degrees.

\subsection{Admissibility hypotheses}
\label{subsec:admissibility-hypotheses}

\begin{definition}[Mealy-effective interface cover]
\label{def:mealy-effective-interface-cover}
An internal functor
\[
    p:\mathbb U\to\mathbb A
\]
is a \textbf{Mealy-effective interface cover} if:
\begin{enumerate}
    \item the relevant slice pullback functor along
    \[
        p_0\times 1_{B_0}:
        U_0\times B_0
        \to
        A_0\times B_0
    \]
    is monadic;
    \item descent is effective for the derivative-domain diagrams obtained
    from \(p_1\), \(p_0\), and descended carriers;
    \item compatible maps into the fixed objects \(B_0\) and \(B_1\) descend;
    \item equality of maps is reflected by pullback along \(p_0\), \(p_1\),
    and the corresponding maps on composable pairs;
    \item these properties are stable under pullback of \(p\).
\end{enumerate}
\end{definition}

\begin{remark}[Componentwise sufficient conditions]
\label{rem:componentwise-sufficient-descent}
A more concrete sufficient criterion may be imposed by requiring:
\begin{enumerate}
    \item effective descent at the levels carrying carriers and primitive
    transition data;
    \item descent at the level carrying composable pairs;
    \item almost descent or conservativity at the level carrying coherence
    equations.
\end{enumerate}
The exact minimal hierarchy depends on the chosen general theorem for descent
of finite-limit models. Definition~\ref{def:mealy-effective-interface-cover}
records only the properties used in the proof below.
\end{remark}

\subsection{Finite-limit presentation of strict Mealy systems}
\label{subsec:finite-limit-presentation-mealy}

For fixed \((\mathbb A,\mathbb B)\), a strict Mealy system consists of:
\begin{enumerate}
    \item a carrier \(P\);
    \item maps
    \[
        p_P:P\to A_0,
        \qquad
        q_P:P\to B_0;
    \]
    \item the pullback object
    \[
        D_P=A_1\times_{t_A,A_0,p_P}P;
    \]
    \item maps
    \[
        \delta_P:D_P\to P,
        \qquad
        \omega_P:D_P\to B_1;
    \]
    \item finite-limit equations expressing typing, units, and
    multiplication.
\end{enumerate}

Thus strict Mealy systems are models of a finite-limit sketch depending on
\((\mathbb A,\mathbb B)\).

\begin{proposition}[Finite-limit definability]
\label{prop:finite-limit-definability-mealy}
The category
\[
    \mathsf{Sys}_{\mathbb B}(\mathbb A)
\]
is equivalent to a category of models of a finite-limit sketch in
\(\mathcal E\).
\end{proposition}

\begin{proof}
All structure objects are obtained from finite products and pullbacks, and all
axioms are equations between morphisms built from those finite-limit
constructions.
\end{proof}

\section{Effective descent for strict Mealy systems}
\label{sec:effective-descent-strict-mealy}

Assume throughout this section that
\[
    p:\mathbb U\to\mathbb A
\]
is a Mealy-effective interface cover.

\subsection{Descent of carriers}
\label{subsec:descent-carriers}

Let
\[
    (P_{\mathbb U},\vartheta_P)
\]
be the carrier part of a system descent datum.

By effective descent in the slice, there exists a carrier
\[
    P\to A_0\times B_0
\]
and an isomorphism
\[
    p^\ast P\cong P_{\mathbb U}
\]
compatible with the descent datum.

\subsection{Descent of derivative domains}
\label{subsec:descent-derivative-domains}

The derivative domain of the global carrier is
\[
    D_P
    :=
    A_1\times_{t_A,A_0,p_P}P.
\]

Pulling it back along \(p\) gives an object canonically isomorphic to
\[
    U_1\times_{t_U,U_0,p_{P_{\mathbb U}}}P_{\mathbb U}.
\]

\begin{lemma}[Derivative domains commute with descent pullback]
\label{lem:derivative-domains-descent}
There is a canonical isomorphism
\[
    p^\ast D_P
    \cong
    D_{P_{\mathbb U}}.
\]
\end{lemma}

\begin{proof}
Both objects are obtained from the same pullback square after applying the
pullback functor associated to \(p_1\) and \(p_0\). The result follows from
pasting pullback squares.
\end{proof}

\subsection{Descent of transition maps}
\label{subsec:descent-transition-maps}

The local transition map
\[
    \delta_{\mathbb U}:
    D_{P_{\mathbb U}}
    \to
    P_{\mathbb U}
\]
is compatible with the overlap isomorphism.

By effective descent for maps between descended objects, it determines a
unique global map
\[
    \delta_P:D_P\to P.
\]

\subsection{Descent of output maps}
\label{subsec:descent-output-maps}

Similarly, the local output map
\[
    \omega_{\mathbb U}:
    D_{P_{\mathbb U}}
    \to
    B_1
\]
is compatible on overlaps and therefore descends to
\[
    \omega_P:D_P\to B_1.
\]

\subsection{Descent of typing equations}
\label{subsec:descent-typing-equations}

The local typing equations are:
\[
    p_{P_{\mathbb U}}\delta_{\mathbb U}
    =
    s_U\pi_{U_1},
\]
\[
    s_B\omega_{\mathbb U}
    =
    q_{P_{\mathbb U}}\delta_{\mathbb U},
\]
\[
    t_B\omega_{\mathbb U}
    =
    q_{P_{\mathbb U}}\pi_{P_{\mathbb U}}.
\]

After pulling back the corresponding candidate global equations along the
cover, they become the local equations. Since pullback along the cover
reflects equality, the global typing equations hold.

\subsection{Descent of unit and multiplication laws}
\label{subsec:descent-unit-multiplication}

The unit and multiplication laws are equations over objects built from:
\begin{enumerate}
    \item the carrier;
    \item the arrow object of \(\mathbb A\);
    \item the object of composable pairs of arrows;
    \item finite pullbacks of these objects.
\end{enumerate}

After pullback along \(p\), the candidate global laws become the local Mealy
laws. Equality reflection gives the global laws.

\begin{theorem}[Effective descent for strict Mealy systems]
\label{thm:effective-descent-strict-mealy}
For every Mealy-effective interface cover
\[
    p:\mathbb U\to\mathbb A,
\]
the comparison functor
\[
    K_p:
    \mathsf{Sys}_{\mathbb B}(\mathbb A)
    \longrightarrow
    \mathsf{Desc}_p(\mathsf{Sys}_{\mathbb B})
\]
is an equivalence.
\end{theorem}

\begin{proof}
\leavevmode
\begin{enumerate}
    \item \textbf{Essential surjectivity.}
    Given a system descent datum, descend its carrier, derivative domain,
    transition, and output maps as above. Equality reflection gives all Mealy
    laws. The resulting global system pulls back to the original local system.

    \item \textbf{Fullness.}
    A compatible local strict homomorphism descends to a map of global
    carriers. Transition and output preservation hold globally because they
    hold after pullback along the cover.

    \item \textbf{Faithfulness.}
    If two global strict homomorphisms have equal pullbacks, equality
    reflection implies that the global maps are equal.
\end{enumerate}
\end{proof}

\subsection{The system topology}
\label{subsec:system-topology}

Let
\[
    \tau_{\amalg}
\]
denote the extensive coverage generated by coproduct decompositions, and let
\[
    \tau_{\mathrm{ed}}
\]
denote the singleton coverage generated by Mealy-effective interface covers.

\begin{definition}[System interface topology]
\label{def:system-interface-topology}
The \textbf{system interface topology} is the topology
\[
    \tau_{\mathrm{sys}}
    :=
    \tau_{\amalg}\vee\tau_{\mathrm{ed}}
\]
generated jointly by extensive covers and Mealy-effective singleton covers.
\end{definition}

\begin{corollary}[Stack of strict Mealy systems]
\label{cor:stack-strict-mealy-systems}
The indexed category
\[
    \mathsf{Sys}_{\mathbb B}
\]
is a stack for
\[
    \tau_{\mathrm{sys}}.
\]
\end{corollary}

\begin{proof}
Extensive descent is
Theorem~\ref{thm:extensive-descent-systems}. Effective singleton descent is
Theorem~\ref{thm:effective-descent-strict-mealy}. The topology is generated by
these covers.
\end{proof}

\section{Residual behaviour under input restriction}
\label{sec:residual-behaviour-input-restriction}

The relationship between input restriction and canonical residual behaviour
requires more care than the corresponding relationship for carriers.

\subsection{Restriction of residual categories}
\label{subsec:restriction-residual-categories}

Let
\[
    F:\mathbb A'\to\mathbb A
\]
be an internal functor and let
\[
    v'\in A'_0.
\]
There is an induced internal functor
\[
    F/v':
    \mathbb A'/v'
    \to
    \mathbb A/F_0v'
\]
sending
\[
    r:u'\to v'
\]
to
\[
    F_1r:F_0u'\to F_0v'.
\]

Precomposition gives a restriction functor
\[
    (F/v')^\ast:
    [\mathbb A/F_0v',\mathbb B]
    \to
    [\mathbb A'/v',\mathbb B].
\]

These functors assemble into a comparison map on canonical behaviour
carriers.

\begin{definition}[Residual restriction comparison]
\label{def:residual-restriction-comparison}
The \textbf{residual restriction comparison} is the map
\[
    \rho_F:
    F^\ast
    \mathbf{Beh}_{\mathbb A,\mathbb B,0}
    \to
    \mathbf{Beh}_{\mathbb A',\mathbb B,0}
\]
whose component over \(v'\) sends
\[
    H:\mathbb A/F_0v'\to\mathbb B
\]
to
\[
    H\circ(F/v').
\]
\end{definition}

\begin{remark}[Restriction need not be essentially surjective]
\label{rem:residual-restriction-not-surjective}
An arbitrary residual behaviour on \(\mathbb A'/v'\) need not extend to one
on \(\mathbb A/F_0v'\). Thus \(\rho_F\) is not automatically an isomorphism.
Additional hypotheses on \(F\) are required for behaviour descent.
\end{remark}

\subsection{Behaviour-compatible covers}
\label{subsec:behaviour-compatible-covers}

\begin{definition}[Behaviour-compatible input functor]
\label{def:behaviour-compatible-input-functor}
An internal functor
\[
    F:\mathbb A'\to\mathbb A
\]
is \textbf{behaviour-compatible relative to \(\mathbb B\)} if the residual
restriction comparison
\[
    \rho_F:
    F^\ast
    \mathbf{Beh}_{\mathbb A,\mathbb B,0}
    \to
    \mathbf{Beh}_{\mathbb A',\mathbb B,0}
\]
is an isomorphism.
\end{definition}

\begin{remark}[Relative nature of behaviour compatibility]
\label{rem:relative-behaviour-compatibility}
Behaviour compatibility may depend on \(\mathbb B\). A stronger
output-independent condition would require each residual functor
\[
    F/v':
    \mathbb A'/v'
    \to
    \mathbb A/F_0v'
\]
to induce an equivalence on all internal functor categories into admissible
outputs.
\end{remark}

\begin{proposition}[Derivative compatibility]
\label{prop:residual-restriction-derivative-compatibility}
For every internal functor
\[
    F:\mathbb A'\to\mathbb A,
\]
the residual restriction comparison commutes with derivatives:
\[
    \rho_F\left(
    \partial_{F_1a'}H
    \right)
    =
    \partial_{a'}\left(\rho_FH\right).
\]
\end{proposition}

\begin{proof}
Both sides are obtained by precomposition with the residual functor induced by
\(a'\) and then with \(F\). The equality follows from functoriality of slice
reindexing.
\end{proof}

\section{Descent of contracts}
\label{sec:descent-contracts}

Assume throughout this section that
\[
    p:\mathbb U\to\mathbb A
\]
is:
\begin{enumerate}
    \item a Mealy-effective interface cover;
    \item behaviour-compatible relative to \(\mathbb B\);
    \item effective for subobject descent on the canonical behaviour carrier.
\end{enumerate}

\subsection{Descent of behavioural predicates}
\label{subsec:descent-behavioural-predicates}

Since
\[
    p^\ast
    \mathbf{Beh}_{\mathbb A,\mathbb B,0}
    \cong
    \mathbf{Beh}_{\mathbb U,\mathbb B,0},
\]
subobject descent gives a comparison
\[
    \mathsf{Pred}_{\mathbb B}(\mathbb A)
    \to
    \mathsf{Desc}_p(\mathsf{Pred}_{\mathbb B}).
\]

\begin{proposition}[Descent of behavioural predicates]
\label{prop:descent-behavioural-predicates}
Under the standing assumptions, the comparison
\[
    \mathsf{Pred}_{\mathbb B}(\mathbb A)
    \longrightarrow
    \mathsf{Desc}_p(\mathsf{Pred}_{\mathbb B})
\]
is an isomorphism of posets.
\end{proposition}

\begin{proof}
This is effective descent for subobjects of the descended canonical behaviour
object.
\end{proof}

\subsection{Locality of derivative invariance}
\label{subsec:locality-derivative-invariance}

Let
\[
    C\hookrightarrow
    \mathbf{Beh}_{\mathbb A,\mathbb B,0}
\]
be a behavioural predicate.

\begin{proposition}[Derivative invariance is local]
\label{prop:derivative-invariance-local}
The predicate \(C\) is derivative-invariant if and only if its pullback
\[
    p^\ast C
\]
is derivative-invariant.
\end{proposition}

\begin{proof}
If \(C\) is derivative-invariant, the result follows from
Proposition~\ref{prop:residual-restriction-derivative-compatibility}.

Conversely, suppose \(p^\ast C\) is derivative-invariant. The global
derivative-closure factorization exists after pulling back along the cover.
Since the cover is effective for the relevant subobject factorization and
reflects inclusion, the derivative of \(C\) factors through \(C\) globally.
\end{proof}

\begin{theorem}[Contract descent]
\label{thm:contract-descent}
Under the standing assumptions, the comparison
\[
    \mathsf{Ctr}_{\mathbb B}(\mathbb A)
    \longrightarrow
    \mathsf{Desc}_p(\mathsf{Ctr}_{\mathbb B})
\]
is an isomorphism of posets.
\end{theorem}

\begin{proof}
By Proposition~\ref{prop:descent-behavioural-predicates}, compatible local
predicates descend uniquely. By
Proposition~\ref{prop:derivative-invariance-local}, the descended predicate is
a residual contract exactly when the local predicates are residual contracts.
\end{proof}

\subsection{The generation--inclusion--safety adjunction under restriction}
\label{subsec:indexed-generation-inclusion-safety}

For each \(\mathbb A\), recall the adjoint triple
\[
    \operatorname{Gen}_{\partial,\mathbb A}
    \dashv
    I_{\partial,\mathbb A}
    \dashv
    \operatorname{Safe}_{\partial,\mathbb A}.
\]

Let
\[
    F:\mathbb A'\to\mathbb A
\]
be behaviour-compatible.

\begin{proposition}[Generated closure and input restriction]
\label{prop:generated-closure-input-restriction}
Assume regular images are stable under the relevant pullbacks. Then there is a
canonical isomorphism
\[
    F^\ast
    \operatorname{Gen}_{\partial,\mathbb A}(C)
    \cong
    \operatorname{Gen}_{\partial,\mathbb A'}
    (F^\ast C).
\]
\end{proposition}

\begin{proof}
Generated closure is the regular image of the derivative action on \(C\).
Derivative compatibility identifies the pulled-back derivative diagram with
the derivative diagram of \(F^\ast C\). Stability of regular images under
pullback gives the result.
\end{proof}

\begin{proposition}[Safety core and input restriction]
\label{prop:safety-core-input-restriction}
Assume the dependent products defining the safety cores satisfy
Beck--Chevalley for the relevant pullback squares. Then there is a canonical
isomorphism
\[
    F^\ast
    \operatorname{Safe}_{\partial,\mathbb A}(C)
    \cong
    \operatorname{Safe}_{\partial,\mathbb A'}
    (F^\ast C).
\]
\end{proposition}

\begin{proof}
The safety core is
\[
    \forall_{\pi_{\mathbf{Beh}}}
    \left(
    \partial^\ast C
    \right).
\]
Derivative compatibility identifies the pulled-back predicate
\(\partial^\ast C\) with the corresponding derivative predicate over
\(\mathbb A'\). Beck--Chevalley for the dependent product gives the claimed
isomorphism.
\end{proof}

\begin{corollary}[Indexed adjoint triple]
\label{cor:indexed-adjoint-triple}
On the subcategory of behaviour-compatible input interfaces satisfying the
stated regular-image and Beck--Chevalley hypotheses, the fibrewise adjoint
triples
\[
    \operatorname{Gen}_{\partial}
    \dashv
    I_{\partial}
    \dashv
    \operatorname{Safe}_{\partial}
\]
assemble into an indexed adjoint triple.
\end{corollary}

\section{Behavioural abstraction under descent}
\label{sec:behavioural-abstraction-descent}

\subsection{Semantic maps and input restriction}
\label{subsec:semantic-maps-input-restriction}

Let
\[
    \Phi:\mathbb A\rightsquigarrow\mathbb B
\]
have terminal semantic map
\[
    \beta_\Phi:
    P\to
    \mathbf{Beh}_{\mathbb A,\mathbb B,0}.
\]

For
\[
    F:\mathbb A'\to\mathbb A,
\]
there are two maps from the pulled-back carrier:
\begin{enumerate}
    \item pull back \(\beta_\Phi\) and then apply residual restriction;
    \item form the terminal semantic map of \(F^\ast\Phi\).
\end{enumerate}

\begin{proposition}[Naturality of residual semantics]
\label{prop:naturality-residual-semantics-input}
The square
\[
\begin{tikzcd}[column sep=large,row sep=large]
    F^\ast P
        \arrow[r, "F^\ast\beta_\Phi"]
        \arrow[d, "\beta_{F^\ast\Phi}"']
    &
    F^\ast
    \mathbf{Beh}_{\mathbb A,\mathbb B,0}
        \arrow[d, "\rho_F"]
    \\
    \mathbf{Beh}_{\mathbb A',\mathbb B,0}
        \arrow[r, equal]
    &
    \mathbf{Beh}_{\mathbb A',\mathbb B,0}
\end{tikzcd}
\]
commutes.
\end{proposition}

\begin{proof}
For a pulled-back state \((v',x)\), the upper-right composite assigns the
residual behaviour of \(x\) and then restricts it along
\[
    \mathbb A'/v'\to\mathbb A/F_0v'.
\]
The lower-left composite computes precisely the residual behaviour of the
restricted system from \((v',x)\). These agree on residual objects and arrows.
\end{proof}

\subsection{Locality of semantic images}
\label{subsec:locality-semantic-images}

Assume that \(\mathcal E\) is regular.

Write
\[
    \operatorname{Beh}_{\mathbb A}(\Phi)
    \hookrightarrow
    \mathbf{Beh}_{\mathbb A,\mathbb B,0}
\]
for the regular image of \(\beta_\Phi\).

\begin{proposition}[Behavioural abstraction and input restriction]
\label{prop:behavioural-abstraction-input-restriction}
Let
\[
    F:\mathbb A'\to\mathbb A
\]
be behaviour-compatible, and assume regular images are stable under the
relevant pullback. Then there is a canonical isomorphism
\[
    F^\ast
    \operatorname{Beh}_{\mathbb A}(\Phi)
    \cong
    \operatorname{Beh}_{\mathbb A'}(F^\ast\Phi).
\]
\end{proposition}

\begin{proof}
By Proposition~\ref{prop:naturality-residual-semantics-input}, the semantic map
of the restricted system is the pullback of the global semantic map followed
by the behaviour-comparison isomorphism. Stability of regular images under
pullback identifies the two semantic images.
\end{proof}

\begin{remark}[Systems interpretation]
\label{rem:black-boxing-local-observation}
The proposition says that black-box abstraction commutes with restriction to
an admissible input-interface patch. One may either:
\begin{enumerate}
    \item abstract the global system and then inspect the local input view; or
    \item restrict the system first and then compute its behaviour.
\end{enumerate}
The resulting local behavioural contract is the same.
\end{remark}

\subsection{Locality of contract satisfaction}
\label{subsec:locality-contract-satisfaction}

Let
\[
    V\hookrightarrow
    \mathbf{Beh}_{\mathbb A,\mathbb B,0}
\]
be a residual behavioural contract.

Recall that
\[
    \Phi\models V
\]
means that
\[
    \beta_\Phi
\]
factors through \(V\), equivalently
\[
    \operatorname{Beh}_{\mathbb A}(\Phi)\leq V.
\]

\begin{theorem}[Locality of contract satisfaction]
\label{thm:locality-contract-satisfaction}
Let
\[
    \{F_i:\mathbb A_i\to\mathbb A\}_{i\in I}
\]
be a covering family which is:
\begin{enumerate}
    \item effective for strict Mealy systems;
    \item effective for residual contracts;
    \item jointly conservative on subobjects of the canonical behaviour
    object;
    \item behaviour-compatible.
\end{enumerate}
Then
\[
    \Phi\models V
\]
if and only if
\[
    F_i^\ast\Phi
    \models
    F_i^\ast V
\]
for every \(i\in I\).
\end{theorem}

\begin{proof}
If \(\Phi\models V\), pullback preserves the factorization, so every local
system models the local contract.

Conversely, suppose every local restriction models the corresponding local
contract. By
Proposition~\ref{prop:behavioural-abstraction-input-restriction},
\[
    F_i^\ast
    \operatorname{Beh}_{\mathbb A}(\Phi)
    \leq
    F_i^\ast V
\]
for every \(i\). Joint conservativity of the covering family implies
\[
    \operatorname{Beh}_{\mathbb A}(\Phi)\leq V.
\]
Hence \(\Phi\models V\).
\end{proof}

\subsection{Abstraction--realization under input restriction}
\label{subsec:abstraction-realization-input-restriction}

Recall the abstraction--realization adjunction
\[
    \operatorname{Beh}_{\mathbb A}
    \dashv
    \iota_{\mathbb A}.
\]

\begin{proposition}[Restriction of canonical realizations]
\label{prop:restriction-canonical-realizations}
Let
\[
    F:\mathbb A'\to\mathbb A
\]
be behaviour-compatible. For every residual contract
\[
    V\in\mathsf{Ctr}_{\mathbb B}(\mathbb A),
\]
there is a canonical isomorphism
\[
    F^\ast\iota_{\mathbb A}(V)
    \cong
    \iota_{\mathbb A'}(F^\ast V).
\]
\end{proposition}

\begin{proof}
The canonical realization of \(V\) is the restriction of the canonical
behaviour Mealy morphism to the derivative-invariant subobject \(V\).
Behaviour compatibility identifies the pulled-back canonical behaviour
machine with the local canonical behaviour machine, while derivative
compatibility identifies the restricted transition-output structures.
\end{proof}

\begin{corollary}[Indexed abstraction--realization adjunction]
\label{cor:indexed-abstraction-realization}
On the admissible input-interface site, the fibrewise adjunctions
\[
    \operatorname{Beh}_{\mathbb A}
    \dashv
    \iota_{\mathbb A}
\]
assemble into an indexed adjunction between the stack of strict Mealy systems
and the stack of residual behavioural contracts.
\end{corollary}

\section{Temporal--interface interchange}
\label{sec:temporal-interface-interchange}

\subsection{Execution categories under input pullback}
\label{subsec:execution-categories-input-pullback}

Let
\[
    F:\mathbb A'\to\mathbb A
\]
and
\[
    \Phi=(P,\delta_P,\omega_P):
    \mathbb A\rightsquigarrow\mathbb B.
\]

The restricted carrier is
\[
    P_F=A'_0\times_{A_0}P.
\]

The execution category of \(F^\ast\Phi\) has:
\[
    \left(
    \int_{\mathbb A'}P_F
    \right)_0
    =
    P_F,
\]
and
\[
    \left(
    \int_{\mathbb A'}P_F
    \right)_1
    =
    A'_1\times_{A'_0}P_F.
\]

There is a canonical internal functor
\[
    E_{F,\Phi}:
    \int_{\mathbb A'}P_F
    \to
    \int_{\mathbb A}P
\]
defined by
\[
    E_{F,\Phi}(v',x)=x
\]
on objects and
\[
    E_{F,\Phi}(a',(v',x))
    =
    (F_1a',x)
\]
on arrows.

\begin{proposition}[Execution restriction functor]
\label{prop:execution-restriction-functor}
The displayed maps define an internal functor
\[
    E_{F,\Phi}:
    \int_{\mathbb A'}F^\ast P
    \to
    \int_{\mathbb A}P.
\]
It commutes with the input and output labelling:
\[
    I_P E_{F,\Phi}
    =
    F^{\mathrm{op}}I_{F^\ast P},
\]
and
\[
    O_P E_{F,\Phi}
    =
    O_{F^\ast P}.
\]
\end{proposition}

\begin{proof}
Source preservation is immediate. Target preservation is the defining
equation
\[
    \delta_{F^\ast P}(a',(v',x))
    =
    \left(
    s_{A'}a',
    \delta_P(F_1a',x)
    \right).
\]
Identity and composition preservation follow from functoriality of \(F\).
The labelling equations follow from the definitions of input pullback.
\end{proof}

\subsection{Trajectory restriction}
\label{subsec:trajectory-restriction-input}

Applying temporal trajectories gives a morphism
\[
    \operatorname{Temp}(E_{F,\Phi}):
    \mathscr T_{F^\ast\Phi}
    \to
    \mathscr T_\Phi.
\]

More precisely, the local trajectory sheaf is obtained by restricting the
input labels of the global trajectory sheaf along \(F\).

\begin{definition}[Interface restriction of trajectory semantics]
\label{def:interface-restriction-trajectory-semantics}
Write
\[
    F^\sharp\mathscr T_\Phi
\]
for the pullback
\[
\begin{tikzcd}[column sep=large]
    F^\sharp\mathscr T_\Phi
        \arrow[r]
        \arrow[d]
    &
    \mathscr T_\Phi
        \arrow[d, "\operatorname{Temp}(I_P)"]
    \\
    \operatorname{Temp}(\mathbb A'^{\mathrm{op}})
        \arrow[r, "\operatorname{Temp}(F^{\mathrm{op}})"']
    &
    \operatorname{Temp}(\mathbb A^{\mathrm{op}}).
\end{tikzcd}
\]
\end{definition}

\begin{theorem}[Trajectory semantics commutes with input restriction]
\label{thm:trajectory-semantics-input-restriction}
There is a canonical isomorphism of interval sheaves
\[
    \mathscr T_{F^\ast\Phi}
    \cong
    F^\sharp\mathscr T_\Phi.
\]
\end{theorem}

\begin{proof}
At temporal length \(0\), both sides are the pulled-back carrier
\[
    A'_0\times_{A_0}P.
\]
At temporal length \(1\), both sides are
\[
    A'_1\times_{A'_0}
    \left(
    A'_0\times_{A_0}P
    \right),
\]
with the same input and output labels.

For general \(n\), both sides are obtained by the same iterated pullback of
the length-one execution object over the length-zero state object. The
canonical lengthwise isomorphisms commute with temporal restriction maps.
\end{proof}

\subsection{The two-variable execution object}
\label{subsec:two-variable-execution-object}

Fix a global input interface \(\mathbb A\) and a system
\[
    \Phi:\mathbb A\rightsquigarrow\mathbb B.
\]

Let
\[
    \mathsf{Patch}(\mathbb A)
\]
be a chosen full subcategory of
\[
    \mathsf{IntCat}(\mathcal E)/\mathbb A
\]
whose objects are admissible input patches
\[
    u:\mathbb U\to\mathbb A.
\]

\begin{definition}[Temporal--interface execution object]
\label{def:temporal-interface-execution-object}
Define
\[
    \mathscr X_\Phi:
    \mathbf{Int}_{\mathbb N}^{\mathrm{op}}
    \times
    \mathsf{Patch}(\mathbb A)^{\mathrm{op}}
    \to
    \mathcal E
\]
by
\[
    \mathscr X_\Phi(n,u)
    :=
    \mathscr T_{u^\ast\Phi}(n).
\]
\end{definition}

For fixed \(u\), the object
\[
    n\longmapsto\mathscr X_\Phi(n,u)
\]
is a discrete interval sheaf.

For fixed \(n\), the object
\[
    u\longmapsto\mathscr X_\Phi(n,u)
\]
is an interface-indexed object built from finite limits of the restricted
system.

\subsection{Interface descent of finite trajectories}
\label{subsec:interface-descent-finite-trajectories}

Let
\[
    p:\mathbb U\to\mathbb A
\]
be a Mealy-effective interface cover.

Since \(n\)-step execution objects are constructed from:
\begin{enumerate}
    \item the descended carrier;
    \item the descended arrow object;
    \item finitely many pullbacks;
\end{enumerate}
they inherit effective descent.

\begin{proposition}[Descent of \(n\)-step executions]
\label{prop:descent-n-step-executions}
For every \(n\in\mathbb N\), the object of \(n\)-step executions satisfies
effective descent along \(p\).
\end{proposition}

\begin{proof}
The case \(n=0\) is descent of carriers. The case \(n=1\) is descent of the
primitive execution object. The general case is obtained by iterated
pullbacks. Effective descent is stable under the finite-limit constructions
used to form these path objects.
\end{proof}

\subsection{Interchange of temporal and interface gluing}
\label{subsec:interchange-temporal-interface-gluing}

Temporal gluing is computed by pullback over a boundary state object:
\[
    \mathscr T_\Phi(m+n)
    \cong
    \mathscr T_\Phi(m)
    \times_{\mathscr T_\Phi(0)}
    \mathscr T_\Phi(n).
\]

Interface descent reconstructs each of these objects from compatible local
data.

\begin{theorem}[Temporal--interface interchange]
\label{thm:temporal-interface-interchange}
Let
\[
    p:\mathbb U\to\mathbb A
\]
be a Mealy-effective interface cover. Then temporal gluing and interface
descent commute.

Equivalently, for every \(m,n\), the canonical comparison between:
\begin{enumerate}
    \item first gluing local \(m\)- and \(n\)-step trajectories temporally and
    then descending in the interface variable;
    \item first descending the local \(m\)- and \(n\)-step trajectory objects
    and then gluing them temporally;
\end{enumerate}
is an isomorphism.
\end{theorem}

\begin{proof}
Temporal gluing is a finite pullback. The objects and maps entering this
pullback satisfy effective descent by
Proposition~\ref{prop:descent-n-step-executions}. Since the descent
reconstruction preserves the relevant finite limits, the pullback computed
before descent agrees canonically with the pullback computed after descent.
\end{proof}

\begin{corollary}[Double locality of execution]
\label{cor:double-locality-execution}
The execution semantics of a strict Mealy system is:
\begin{enumerate}
    \item a sheaf in the temporal interval variable;
    \item compatible with effective descent in the input-interface variable.
\end{enumerate}
\end{corollary}

\begin{remark}[Systems interpretation]
\label{rem:systems-interpretation-interchange}
A global finite execution may be reconstructed in either order:
\begin{enumerate}
    \item glue temporal segments within each local interface patch and then
    integrate the patches;
    \item integrate interface-local execution data and then glue the resulting
    global temporal segments.
\end{enumerate}
The interchange theorem says that these procedures agree.
\end{remark}

\section{Systems-theoretic examples}
\label{sec:systems-theoretic-descent-examples}

\subsection{Disconnected operating modes}
\label{subsec:example-disconnected-operating-modes}

Let
\[
    \mathbb A=\mathbb A_1\amalg\mathbb A_2.
\]
A system
\[
    \Phi:\mathbb A\rightsquigarrow\mathbb B
\]
is equivalently a pair
\[
    \Phi_1:\mathbb A_1\rightsquigarrow\mathbb B,
    \qquad
    \Phi_2:\mathbb A_2\rightsquigarrow\mathbb B.
\]

This models two operating modes with no input transition between them. A
residual contract over \(\mathbb A\) is exactly a pair of contracts over the
two modes.

\subsection{Overlapping sensor views}
\label{subsec:example-overlapping-sensor-views}

Suppose a global input interface records variables
\[
    x,\ y,\ z.
\]
Let
\[
    \mathbb U_{xy}\to\mathbb A
\]
be the patch exposing \(x,y\), and let
\[
    \mathbb U_{yz}\to\mathbb A
\]
be the patch exposing \(y,z\).

Their overlap exposes the shared variable \(y\).

A system descent datum consists of:
\begin{enumerate}
    \item a local system on the \(x,y\)-view;
    \item a local system on the \(y,z\)-view;
    \item an isomorphism between their restrictions to the \(y\)-view;
    \item coherence of this isomorphism on higher overlaps.
\end{enumerate}

Effective descent says that such compatible local sensor models arise from a
global system on \(x,y,z\).

\subsection{Distributed controller integration}
\label{subsec:example-distributed-controller-integration}

Let two local controllers observe overlapping disturbance interfaces:
\[
    \mathbb U_1\to\mathbb A,
    \qquad
    \mathbb U_2\to\mathbb A.
\]
Each controller determines a local Mealy system
\[
    \Phi_i:\mathbb U_i\rightsquigarrow\mathbb B.
\]

The overlap comparison requires the two controllers to agree on:
\begin{enumerate}
    \item common disturbance inputs;
    \item common state identifications;
    \item emitted control actions on shared situations.
\end{enumerate}

A global descended system represents a coordinated controller. Failure of the
cocycle condition represents a coordination obstruction.

\subsection{Pairwise agreement without triple coherence}
\label{subsec:example-pairwise-without-triple-coherence}

Consider three patches
\[
    \mathbb U_1,\mathbb U_2,\mathbb U_3
\]
with pairwise overlap isomorphisms
\[
    \vartheta_{12},
    \qquad
    \vartheta_{23},
    \qquad
    \vartheta_{13}.
\]

Suppose
\[
    \vartheta_{23}\vartheta_{12}
    \neq
    \vartheta_{13}
\]
on the triple overlap.

Then the local systems are pairwise compatible but do not form a descent
datum. No global system can induce these overlap identifications.

This is a higher coherence obstruction rather than a failure visible on any
single pair of patches.

\subsection{A global contract without a global implementation}
\label{subsec:example-global-contract-no-implementation}

Suppose compatible local contracts
\[
    V_i
\]
descend to a global residual contract
\[
    V.
\]
It may nevertheless happen that the local systems chosen to realize the
\(V_i\) do not admit compatible overlap identifications.

Thus:
\[
    \text{contract descent}
\]
does not imply
\[
    \text{implementation descent}.
\]

The distinction separates global specification consistency from global
realizability.

\subsection{A global implementation with inconsistent local specifications}
\label{subsec:example-global-implementation-inconsistent-specifications}

Conversely, a global system \(\Phi\) may exist while independently chosen
local contracts fail to agree on overlaps. The obstruction is then
specification inconsistency rather than implementation inconsistency.

\subsection{Genuinely local state data}
\label{subsec:example-genuinely-local-state-data}

Let
\[
    \mathcal E=\operatorname{Sh}(\mathcal C,J)
\]
be a nontrivial Grothendieck topos.

A carrier object \(P\) may have local sections over a covering family
\[
    U_i\to U
\]
without possessing a global section over \(U\). Nevertheless, \(P\) is a
perfectly valid internal state object.

A Mealy transition may therefore evolve locally represented states which do
not admit a single global point. Effective descent applies to the internal
object and its structure maps rather than to a chosen set of global states.

\subsection{Temporal--interface interchange example}
\label{subsec:example-temporal-interface-interchange}

Let a two-step execution be covered temporally by two one-step segments and
interface-wise by two overlapping patches.

One may:
\begin{enumerate}
    \item glue the two temporal segments in each patch, obtaining two local
    two-step executions, and then descend them;
    \item descend each family of one-step executions first, obtaining two
    global one-step executions, and then concatenate them.
\end{enumerate}

Theorem~\ref{thm:temporal-interface-interchange} identifies the resulting
global two-step executions.

\section{Comparison with interval-sheaf abstract machines}
\label{sec:comparison-interval-sheaf-machines}

\subsection{Common span architecture}
\label{subsec:common-span-architecture}

Both the present framework and interval-sheaf approaches to systems associate
to a machine a span
\[
    \text{input behaviour}
    \longleftarrow
    \text{system behaviour}
    \longrightarrow
    \text{output behaviour}.
\]

In the present setting, the derived temporal span is
\[
    \operatorname{Temp}(\mathbb A^{\mathrm{op}})
    \longleftarrow
    \mathscr T_\Phi
    \longrightarrow
    \operatorname{Temp}(\mathbb B^{\mathrm{op}}).
\]

Composition is implemented by pullback compatibility over the shared
intermediate trajectory object.

\subsection{Difference in primitive structure}
\label{subsec:difference-primitive-structure}

In an interval-sheaf abstract machine, temporally extended behaviour is taken
as primitive.

In the present framework, the primitive data are:
\begin{enumerate}
    \item a persistent state carrier;
    \item typed input processing;
    \item transition;
    \item output comparison.
\end{enumerate}

The temporal trajectory sheaf is derived from the resulting execution
category.

\subsection{Segal refinement}
\label{subsec:segal-refinement-comparison}

Because the temporal sheaf arises from an internal category, it has a
canonical Segal enhancement.

The interval sheaf records finite trajectories and temporal restriction. The
Segal enhancement additionally records:
\begin{enumerate}
    \item identity executions;
    \item contraction of adjacent executions by composition;
    \item coherence of iterated composition.
\end{enumerate}

\subsection{Interface descent as a second locality principle}
\label{subsec:interface-descent-second-locality}

The temporal sheaf accounts for locality in time. The stack
\[
    \mathsf{Sys}_{\mathbb B}
\]
accounts for locality in the input interface.

The combined semantics therefore has two local-to-global directions:
\[
    \text{temporal gluing}
    \qquad\text{and}\qquad
    \text{interface reconstruction}.
\]

\subsection{Limitations}
\label{subsec:limitations-temporal-interface}

The constructions of this chapter do not by themselves provide:
\begin{enumerate}
    \item continuous-time trajectory sheaves;
    \item hybrid execution semantics;
    \item a classification of all interval-sheaf machines by Mealy
    morphisms;
    \item effective descent for general Mealy simulations;
    \item simultaneous descent in both input and output interfaces;
    \item temporal fixed-point logic;
    \item cubical or pomset execution semantics.
\end{enumerate}

These require additional structure.

\subsection{Relation to later chapters}
\label{subsec:relation-later-chapters}

The later relative program chapter evaluates Mealy output comparisons in
downstream actions and constructs internal program categories. Its diamonds
and boxes may be studied under input-interface restriction using the
Beck--Chevalley methods established there.

The polynomial-response chapter curries labelled one-step execution into total
response profiles. The present descent theory suggests studying those
polynomial coalgebras as objects indexed over input interfaces.

The concurrency chapters may replace the linear interval site by execution
shapes retaining independence, such as cubical or partially ordered shapes.

The traced chapters may investigate whether feedback closure commutes with
input-interface restriction and descent.

\section{Chapter summary}
\label{sec:temporal-interface-chapter-summary}

\begin{theorem}[Temporal trajectories and interface descent]
\label{thm:temporal-trajectories-interface-descent-summary}
Under the hypotheses stated in the preceding sections, the following hold.

\begin{enumerate}
    \item Internal graphs in \(\mathcal E\) are equivalent to
    \(\mathcal E\)-valued discrete interval sheaves.

    \item Every internal category determines a discrete interval sheaf of
    paths together with a canonical Segal enhancement given by its nerve.

    \item Every Mealy morphism
    \[
        \Phi:\mathbb A\rightsquigarrow\mathbb B
    \]
    determines a temporal execution sheaf
    \[
        \mathscr T_\Phi
        =
        \operatorname{Temp}
        \left(
        \int_{\mathbb A}P
        \right)
    \]
    and a Segal-enhanced span
    \[
        N(\mathbb A^{\mathrm{op}})
        \longleftarrow
        N\left(\int_{\mathbb A}P\right)
        \longrightarrow
        N(\mathbb B^{\mathrm{op}}).
    \]

    \item Temporal execution semantics is functorial in strict semantic
    homomorphisms.

    \item Temporal execution semantics preserves composition of Mealy
    morphisms by pullback over the shared intermediate trajectory sheaf.

    \item For fixed \(\mathbb B\), strict Mealy systems form an indexed
    category
    \[
        \mathsf{Sys}_{\mathbb B}:
        \mathsf{IntCat}(\mathcal E)^{\mathrm{op}}
        \to
        \mathbf{Cat}.
    \]

    \item The indexed category of strict Mealy systems satisfies extensive
    descent over coproduct decompositions of input interfaces.

    \item For every Mealy-effective interface cover
    \[
        p:\mathbb U\to\mathbb A,
    \]
    the comparison
    \[
        \mathsf{Sys}_{\mathbb B}(\mathbb A)
        \to
        \mathsf{Desc}_p(\mathsf{Sys}_{\mathbb B})
    \]
    is an equivalence.

    \item The system doctrine is therefore a stack for the topology
    \[
        \tau_{\mathrm{sys}}
        =
        \tau_{\amalg}\vee\tau_{\mathrm{ed}}.
    \]

    \item Behaviour-compatible input functors induce restriction maps on
    canonical residual behaviours which commute with derivatives.

    \item Residual behavioural contracts satisfy descent for covers which are
    both system-effective and behaviour-compatible.

    \item Under the stated regular-image and Beck--Chevalley hypotheses, the
    adjoint triple
    \[
        \operatorname{Gen}_{\partial}
        \dashv
        I_{\partial}
        \dashv
        \operatorname{Safe}_{\partial}
    \]
    is compatible with input restriction.

    \item Behavioural abstraction commutes with admissible input restriction:
    \[
        F^\ast
        \operatorname{Beh}_{\mathbb A}(\Phi)
        \cong
        \operatorname{Beh}_{\mathbb A'}(F^\ast\Phi).
    \]

    \item Contract satisfaction is local for jointly conservative admissible
    covers:
    \[
        \Phi\models V
        \quad\Longleftrightarrow\quad
        F_i^\ast\Phi\models F_i^\ast V
        \text{ for every }i.
    \]

    \item The abstraction--realization adjunction
    \[
        \operatorname{Beh}\dashv\iota
    \]
    is compatible with admissible input restriction.

    \item Temporal trajectory formation commutes with input restriction:
    \[
        \mathscr T_{F^\ast\Phi}
        \cong
        F^\sharp\mathscr T_\Phi.
    \]

    \item Finite execution objects satisfy effective interface descent.

    \item Temporal gluing and interface descent commute.
\end{enumerate}
\end{theorem}

\begin{proof}
\leavevmode
\begin{enumerate}
    \item Items (1)--(2) are
    Theorem~\ref{thm:internal-graph-interval-sheaf-equivalence} and
    Theorem~\ref{thm:segal-enhancement}.

    \item Items (3)--(5) are
    Definition~\ref{def:temporal-execution-sheaf},
    Proposition~\ref{prop:functoriality-temporal-execution}, and
    Theorem~\ref{thm:temporal-composition}.

    \item Items (6)--(9) are
    Definition~\ref{def:input-indexed-system-category},
    Theorem~\ref{thm:extensive-descent-systems},
    Theorem~\ref{thm:effective-descent-strict-mealy}, and
    Corollary~\ref{cor:stack-strict-mealy-systems}.

    \item Items (10)--(12) are
    Proposition~\ref{prop:residual-restriction-derivative-compatibility},
    Theorem~\ref{thm:contract-descent}, and
    Corollary~\ref{cor:indexed-adjoint-triple}.

    \item Items (13)--(15) are
    Proposition~\ref{prop:behavioural-abstraction-input-restriction},
    Theorem~\ref{thm:locality-contract-satisfaction}, and
    Corollary~\ref{cor:indexed-abstraction-realization}.

    \item Items (16)--(18) are
    Theorem~\ref{thm:trajectory-semantics-input-restriction},
    Proposition~\ref{prop:descent-n-step-executions}, and
    Theorem~\ref{thm:temporal-interface-interchange}.
\end{enumerate}
\end{proof}

\begin{remark}[Final conceptual summary]
\label{rem:temporal-interface-final-summary}
A categorical automaton is local in two distinct senses.

Its executions are local in time: finite trajectories restrict to contiguous
subtrajectories and glue uniquely along matching temporal boundaries.

Its architecture is local in the input interface: under admissible descent
hypotheses, compatible systems on local interface patches reconstruct a
global system.

Residual behaviours, behavioural contracts, and black-box abstractions respect
this second locality whenever the input covers preserve the residual
semantics. Temporal trajectory formation and interface reconstruction then
commute. Thus the same global execution may be obtained either by first
assembling its temporal pieces and then integrating its interface-local
descriptions, or by reversing those two constructions.
\end{remark}